% Level 1, Week 8 -- Loops in Python   (input by lesson-plans.tex)
\lesson{8}{Loops in Python}{A \texttt{for} loop plus pushed-in lines is Scratch's repeat block, written in words.}
\lessoninfo{Concepts: for/range, indentation, loops (R), angles 360/n (R), nested loop (light)}{Tools: Thonny + turtle}{CSTA: 1B-AP-10, 1B-AP-15, 1B-AP-16}

\begin{objectives}
  \item The student can write a \texttt{for i in range(n):} loop and say how many times it repeats (L1.G2, L1.G9; L1.W08.1)
  \item The student can use indentation to show which lines are inside a loop (L1.G2; L1.W08.2)
  \item The student can draw any regular polygon using the rule \enquote{turn = 360 ÷ sides}, and a pattern with a loop inside a loop (L1.G2; L1.W08.3)
  \item The student can find and fix an indentation bug and an off-by-one bug, and ask for help with \enquote{I expected … but …} (L1.G6, L1.G7; L1.W08.4)
\end{objectives}

\begin{materials}
  \item Print: workbook chapter 8 and homework 8 (one per student).
  \item 6 sheets of paper for the line-up cards. Write one line on each, big: \texttt{for i in range(3):}, \texttt{clap}, \texttt{stomp}, \texttt{wave}, \texttt{sit down}, \texttt{say "done!"}.
  \item On each Mac: Thonny open; Finder at \texttt{kids-coding-curriculum/1-introduction/exercises/week08/}.
  \item A timer, sticky notes. Prep time: about 8 minutes.
\end{materials}

\needspace{20\baselineskip}
\section{At a glance}
Times are a guide: if the group is deep in a good build, let it run and shorten the demo.
\begin{timeplan}
  \timeslot{0--8}{Warm-up}{Unplugged \enquote{Indent the line-up}: a human loop.}
  \timeslot{8--18}{Concept}{\texttt{for} and \texttt{range}; the colon and the 4-space push; turn = 360 ÷ sides.}
  \timeslot{18--28}{Guided demo}{Turn last week's 8-line square into a 3-line loop (Exercise 8.1).}
  \timeslot{28--52}{Hands-on build}{Pairs: Exercise 8.1 (3 min), 8.2 Polygon Party (10 min), 8.3 Square Flower (8 min). Swap driver every 8 minutes; last 3 minutes free-build.}
  \timeslot{52--60}{Wrap-up}{Asking for help well, show-and-tell, exit ticket, hand out homework 8.}
\end{timeplan}

\section{Warm-up (0--8 min)}
\begin{unplugged}[Indent the line-up]
Give one card to each student (with fewer students, lay the cards on the table in a column). Arrange them in this order and \textbf{push} the \texttt{clap}, \texttt{stomp} and \texttt{wave} cards one step to the right --- they are \enquote{inside} the loop:
\begin{pycode}
for i in range(3):
    clap
    stomp
    wave
sit down
say "done!"
\end{pycode}
You are the computer: walk down the line and do each card. At \texttt{wave}, go back up to the loop and do the pushed-in cards again, 3 times in total. Then move the \texttt{wave} card back \textbf{out} (not pushed in) and run it again: now \texttt{wave} happens once. Let a student be the computer next.
\end{unplugged}
\begin{teacherscript}
\enquote{Did you see? Moving a card in or out by one step changed what the computer did. In Python, that step is 4 spaces at the start of a line. Spaces matter!}
\end{teacherscript}

\section{Concept (8--18 min)}
\begin{teacherscript}
\enquote{In Scratch, the repeat block hugged the blocks inside it, like a C. Python has no C-shape, so it uses two clues instead: a \textbf{colon} at the end of the \texttt{for} line, and \textbf{4 spaces} in front of every line that is inside. The lines pushed in are the hug.}
\end{teacherscript}
\begin{analogy}[A loop is a playlist on repeat]
\texttt{for i in range(3):} is like pressing \enquote{repeat 3 times} on a playlist. The songs in the playlist are the pushed-in lines. When the last song ends, the playlist starts again --- until it has played 3 times. Then the music player carries on with whatever is next.
\end{analogy}
Show in the Thonny \textbf{Shell} (bottom part), typing one line at a time:
\begin{pycode}
for i in range(4):
    print(i)
\end{pycode}
After the second line press \keys{Return} twice. It prints 0, 1, 2, 3. Say: \enquote{\texttt{range(4)} gives 4 numbers, but Python starts counting at 0. Four numbers --- four turns of the loop.}
Two Thonny helpers: after you type a colon and press \keys{Return}, Thonny pushes the next line in for you. To come back \textbf{out} of the loop, press \keys{Delete} once. \keys{Tab} also adds 4 spaces.
Finally, recall the angle rule from Scratch: \enquote{To get all the way round, the turtle turns 360 degrees in total. A shape with 6 sides turns 360 ÷ 6 = 60 at each corner.}

\section{Guided demo (18--28 min)}
\begin{workedexample}[From 8 lines to 3]
\begin{enumerate}[leftmargin=1.4em,itemsep=1pt]
  \item Open last week's \texttt{week07/ex-first-square/solution.py} (or a student's finished square). Ask: \enquote{Which two lines repeat?} (\texttt{t.forward(100)} and \texttt{t.left(90)}.) \enquote{How many times?} (4.)
  \item Open \texttt{week08/ex-loop-square/starter.py} and run it. A straight line! Ask: \enquote{Why?} (Nothing in the loop turns the turtle.)
  \item Click at the end of \texttt{t.forward(100)}, press \keys{Return} --- Thonny lines you up --- and type \texttt{t.left(90)}. Run. A square from 3 lines of code.
  \item Change \texttt{4} to \texttt{5} and \texttt{90} to \texttt{72}: a pentagon. Ask: \enquote{What about 8 sides?} (45.) Let a student call out numbers.
\end{enumerate}
\end{workedexample}
\pyfile{../exercises/week08/ex-loop-square/solution.py}

\section{Hands-on build (28--52 min)}
\begin{pairroles}[Driver and navigator, with a counting job]
Driver types; navigator reads the README steps and \textbf{counts}: \enquote{How many times will this loop run?} before every Run. Swap at the 8-minute timer.
With 1 student: they drive, you navigate and count. With 3: driver, navigator and checker (runs \texttt{check\_my\_work.py}). With 4: two pairs.
\end{pairroles}
\begin{enumerate}[leftmargin=1.4em,itemsep=2pt]
  \item \textbf{Exercise 8.1 Loop Square} (\texttt{week08/ex-loop-square}): repeat the demo on their own Mac, then run the checker.
  \item \textbf{Exercise 8.2 Polygon Party} (\texttt{week08/ex-polygon-party}): a square and a hexagon, each with its own loop.
  \item \textbf{Exercise 8.3 Square Flower} (\texttt{week08/ex-square-flower}): put the square loop inside a color loop. The inner loop needs 8 spaces.
\end{enumerate}
\begin{bughunt}[Spaces and counting]
Show these two programs on your Mac (type them into a new file) and ask: \enquote{What will each one draw? Why?}
\begin{pycode}[Bug A]
for i in range(3):
    t.forward(100)
    t.left(90)
\end{pycode}
\begin{pycode}[Bug B]
for i in range(4):
    t.forward(100)
t.left(90)
\end{pycode}
Bug A draws only 3 sides (\enquote{off by one}: count the turns of the loop). Bug B draws one long line of 400, then turns once: \texttt{t.left(90)} is not pushed in, so it is outside the loop. Homework 8 has the same two bugs to fix on their own.
\end{bughunt}
\begin{freebuild}[Loop art (last 3 minutes)]
New file \texttt{my-python/loop\_art.py}. Try this magic spiral, then change the numbers:
\begin{pycode}
for i in range(60):
    t.forward(i * 3)
    t.left(91)
\end{pycode}
(Add \texttt{import turtle}, \texttt{t = turtle.Turtle()} at the top and \texttt{turtle.done()} at the end.) What happens with \texttt{left(121)} or \texttt{left(145)}?
\end{freebuild}

\section{Wrap-up (52--60 min)}
\begin{safety}[Asking for help well]
\enquote{Programmers ask for help all the time. Good questions get fast answers.} Teach the help sentence (it is in the workbook too):
\begin{center}\textbf{I expected \rule{2cm}{0.4pt}, but \rule{2cm}{0.4pt} happened. I tried \rule{2cm}{0.4pt}.}\end{center}
Agree: \enquote{Three before me}: read the red message, check the workbook, ask your partner --- then ask the grown-up with the help sentence. Online, the same rule keeps you safe: describe the \textbf{problem}, never share passwords or personal details, and ask a grown-up before posting anywhere.
\end{safety}
Show-and-tell: each pair shows their Square Flower or spiral (30 seconds each).
\begin{formativecheck}[Exit ticket]
\enquote{How many times does it say Hi? How many times Bye?}
\begin{pycode}
for i in range(3):
    print("Hi")
print("Bye")
\end{pycode}
Answer: Hi 3 times, Bye once (\texttt{print("Bye")} is not pushed in). Look for: students who say \enquote{Bye 3 times} need another look at the line-up cards.
\end{formativecheck}
\begin{reflection}
\begin{itemize}[leftmargin=1.2em,itemsep=1pt]
  \item How is \texttt{for i in range(4):} like Scratch's repeat block? How is it different?
  \item Finish the help sentence about a bug you met today.
\end{itemize}
Hand out homework 8 (Loop Fixer).
\end{reflection}

\section{Notes for the teacher}
\begin{commonmistakes}
  \item \textbf{Forgetting the colon} after \texttt{range(4)}: Thonny says \texttt{SyntaxError} or \enquote{expected ':'}. Fix: every \texttt{for} line ends with \texttt{:}.
  \item \textbf{Mixed indentation} (3 spaces on one line, 4 on the next): \texttt{IndentationError}. Fix: select the lines and press \keys{Tab} to push them in, or \keys{Shift}+\keys{Tab} to pull them out.
  \item \textbf{Thinking \texttt{range(4)} counts 1 to 4.} It gives 0, 1, 2, 3. Still 4 times --- that is what matters today.
  \item \textbf{Wrong angle} (using the inside angle, e.g. 120 for a hexagon). Fix: \enquote{How far does the turtle turn? All the way round is 360.}
  \item \textbf{Inner loop not pushed in far enough} in Exercise 8.3: it needs 8 spaces (two pushes).
\end{commonmistakes}
\begin{differentiation}
\supportpath Do Exercise 8.2 only, with the angle written on a sticky note next to the screen (3 → 120, 4 → 90, 6 → 60). Use the line-up cards again if indentation is confusing.
\extensionpath In 8.3, try 12 squares with \texttt{for i in range(12):} and \texttt{t.left(30)}, or triangles instead of squares. Then try a 36-sided shape: it looks like a circle!
\end{differentiation}
\begin{stretchchallenge}[Star power]
Draw a 5-pointed star: 5 lines of 150, turning \texttt{right(144)}. Why 144? (The turtle goes round \textbf{twice}: 5 × 144 = 720.) Can you make a 7-pointed star?
\end{stretchchallenge}
\begin{teachernote}[If you are short on time]
Do the warm-up with 3 cards only (\texttt{for}, \texttt{clap}, \texttt{sit down}). Skip Exercise 8.3 in class --- it makes a lovely start for next week or Buffer B2. Keep the bug hunt: indentation is the big idea today.
\end{teachernote}
